% BHU PYQ -- Manually transcribed & error-corrected
% Paper: MTB-201 : Calculus-II (New Course)
% Exam: B.Sc. (Hons.) Semester II Examination, 2013-14

\documentclass[11pt,a4paper]{article}
\usepackage[a4paper,top=2.5cm,bottom=2.5cm,left=2.8cm,right=2.8cm,headheight=14pt]{geometry}
\usepackage{amsmath,amssymb}
\usepackage[T1]{fontenc}\usepackage[utf8]{inputenc}
\usepackage{lmodern,microtype,fancyhdr,enumitem,hyperref,lastpage,parskip}

\hypersetup{colorlinks=true,urlcolor=black,linkcolor=black,
  pdftitle={MTB-201 Calculus-II -- BHU B.Sc. Sem II 2013-14}}
\pagestyle{fancy}\fancyhf{}
\renewcommand{\headrulewidth}{0.4pt}\renewcommand{\footrulewidth}{0.4pt}
\fancyhead[L]{\small   \textit{Banaras Hindu University}}
\fancyhead[R]{\small   \textit{MTB-201: Calculus-II (New Course)}}
\fancyfoot[C]{\small Page  \thepage\ of \pageref{LastPage}}


\newlist{parts}{enumerate}{2}
\setlist[parts,1]{label=(\alph*),leftmargin=*,itemsep=5pt,topsep=3pt}
\setlist[parts,2]{label=(\roman*),leftmargin=*,itemsep=3pt,topsep=2pt}

\newcommand{\pts}[1]{\hfill   \textit{[#1]}}


\begin{document}

\begin{center}
  {\large\bfseries BANARAS HINDU UNIVERSITY}\\[2pt]
  {
\normalsize B.Sc.\ (Hons.) --- Semester II Examination, 2013-14}\\[6pt]
  \rule{0.8\linewidth}{0.5pt}\\[6pt]
  {\large\bfseries Paper No.\ MTB-201: Calculus-II  extnormal{(New Course)}}\\[4pt]
  {
\normalsize Time: 3 Hours \hfill Full Marks: 70}\\[2pt]
  {\small Code:\ BS/Sem\,II/2}
\end{center}
\rule{\linewidth}{0.4pt}
\smallskip



\noindent   \textit{Instructions: Answer   \textbf{five} questions including Question~1 which is compulsory. Figures in right-hand margin indicate marks.}
\bigskip




\noindent   \textbf{Question 1.}   \textit{(Compulsory.)}

\begin{parts}
  \item Give the definitions of the limit and continuity of $f:\mathbb{R}^2 o\mathbb{R}$ at a point $(a,b)$.\pts{1+1}
  \item Give the definitions of the derivative and partial derivatives of $f:\mathbb{R}^2 o\mathbb{R}$ at $(a,b)$.\pts{1+1}
  \item If $f:\mathbb{R}^2 o\mathbb{R}$ is differentiable at $(a,b)$, prove it is also continuous at $(a,b)$.\pts{2}
  \item Let $f(x,y)=2x^2+3y^2$. Using the $\varepsilon$-$\delta$ definition, prove $\lim_{(x,y) o(1,2)}f(x,y)=14$.\pts{2}
  \item Find $\partial f/\partial y$ for $f(x,y)=\cos y$ at $(a,b)$ using the limit definition.\pts{2}
  \item Let $u:\mathbb{R}^2 o\mathbb{R}$ be continuously differentiable and homogeneous of degree $k$. Prove $x u_x+y u_y=ku$.\pts{2}
  \item If $x=r\cos \theta$, $y=r\sin \theta$, show $r\dfrac{\partial z}{\partial r}=x\dfrac{\partial z}{\partial x}+y\dfrac{\partial z}{\partial y}$ and $\dfrac{\partial z}{\partial \theta}=x\dfrac{\partial z}{\partial y}-y\dfrac{\partial z}{\partial x}$.\pts{1+1}
  \item Evaluate $\displaystyle\int_0^{\pi/2}\!\int_0^{\pi/2}\sin(x+y)\,dx\,dy$.\pts{2}
  \item Compute the volume under $z=1-x-y$ over $D=[0,1] \times[0,1]$.\pts{2}
  \item Change the order of integration: $\displaystyle\int_0^2\!\int_0^{2x}f(x,y)\,dy\,dx$.\pts{2}
  \item Using polar coordinates, find the area of the disc $x^2+y^2\leq a^2$.\pts{2}
\end{parts}
\medskip\rule{\linewidth}{0.2pt}

\medskip


\noindent   \textbf{Question 2.}

\begin{parts}
  \item If the limit of $f$ at $(a,b)$ exists, show it is unique. For $f(x,y)=\dfrac{x^3y}{(x^2+y^2)^2}$, show $\lim_{(x,y) o(0,0)}f$ does not exist.\pts{3+3}
  \item Let $f:\mathbb{R}^2 o\mathbb{R}$ be defined by $f(x,y)=\dfrac{x^2y}{x^4+y^2}$ for $(x,y)
eq(0,0)$ and $f(0,0)=0$. Show $f$ is not continuous at $(0,0)$, but both partial derivatives are defined on all of $\mathbb{R}^2$.\pts{2+4}
\end{parts}
\medskip\rule{\linewidth}{0.2pt}

\medskip


\noindent   \textbf{Question 3.}

\begin{parts}
  \item Define $f(x,y)=\dfrac{xy(x^2-y^2)}{x^2+y^2}$ for $(x,y)
eq(0,0)$ and $f(0,0)=0$. Show $f_{yx}(0,0)
eq f_{xy}(0,0)$.\pts{6}
  \item Let $f(x,y)=x^2y+\sin(1/x)$ for $x
eq0$ and $0$ for $x=0$. Find $f_x$ and $f_y$. Prove $f_x$ is not continuous on the $y$-axis but $f_y$ is continuous everywhere.\pts{6}
\end{parts}
\medskip\rule{\linewidth}{0.2pt}

\medskip


\noindent   \textbf{Question 4.}

\begin{parts}
  \item Let $V$ satisfy the Laplace equation $V_{xx}+V_{yy}=0$. With $x=r\cos \theta$, $y=r\sin \theta$, show $V_{rr}+\dfrac{1}{r}V_r+\dfrac{1}{r^2}V_{ \theta \theta}=0$.\pts{6}
  \item Find the maximum and minimum of $f(x,y)=xy$ on the unit circle $x^2+y^2=1$.\pts{6}
\end{parts}
\medskip\rule{\linewidth}{0.2pt}

\medskip


\noindent   \textbf{Question 5.}

\begin{parts}
  \item Let $g(r, \theta,\phi)=(r\cos \theta\sin\phi,\,r\sin \theta\sin\phi,\,r\cos\phi)$. Compute the Jacobian matrix of $g$ and find the determinant.\pts{6}
  \item State and prove Taylor's theorem for $f(x,y)$.\pts{6}
\end{parts}
\medskip\rule{\linewidth}{0.2pt}

\medskip


\noindent   \textbf{Question 6.}

\begin{parts}
  \item For $f(x,y)=e^{x^2}$, evaluate $\displaystyle\iint_D f\,dx\,dy$ over $D=\{(x,y):0<x<1,\,0\leq y<2x\}$.\pts{6}
  \item Change the order and evaluate $\displaystyle\int_0^a\!\int_0^{\sqrt{a^2-x^2}}\sqrt{x^2+y^2}\,dy\,dx$.\pts{6}
\end{parts}
\medskip\rule{\linewidth}{0.2pt}

\medskip


\noindent   \textbf{Question 7.}

\begin{parts}
  \item For the simplex $V=\{(x,y,z):x\geq0,y\geq0,z\geq0,x+y+z\leq1\}$, prove $\displaystyle\iiint_V x^{\ell-1}y^{m-1}z^{n-1}\,dV=\dfrac{\Gamma(\ell)\Gamma(m)\Gamma(n)}{\Gamma(\ell+m+n+1)}$.\pts{6}
  \item Over the region bounded by the first octant below $\left(\dfrac{x}{a}\right)^p+\left(\dfrac{y}{b}\right)^q+\left(\dfrac{z}{c}\right)^r=1$, show $\displaystyle\iiint x^{\ell-1}y^{m-1}z^{n-1}\,dV=\dfrac{a^\ell b^m c^n\,\Gamma(\ell/p)\Gamma(m/q)\Gamma(n/r)}{pqr\,\Gamma(\ell/p+m/q+n/r+1)}$.\pts{6}
\end{parts}

\vfill\rule{\linewidth}{0.4pt}

\begin{center}\small
\href{https://bhu-exam.vercel.app/}{Click here to visit our website}\\[4pt]
\href{https://me-aryan.vercel.app/}{Click here to visit our contributor}\\[4pt]
\href{https://quashan-me.vercel.app/}{Click here to visit our contributor}
\end{center}
\end{document}
